% f6_terms.tex -- Figure 6: the two terms of eq. (15) in the same whitened coordinates as Figure 4 (natural width about 17.3 cm).
%   (a) d^2: equal on circles about mu (navy dashed); grows in every direction.
%   (b) l: equal on straight lines perpendicular to mu - mu_o; l is the position along that direction times Delta (l = Delta * horizontal
%       coordinate), so it is ONE number, read off a number line.  ID: mean 0, standard deviation about Delta; re-oriented: mean -Delta^2.
%   Same schematic clouds and test images as Figure 4; Delta = 1.768 (CIFAR-10).
\begin{tikzpicture}[>={Stealth[length=1.8mm, width=1.4mm]}, line cap=round, line join=round, font=\footnotesize,
  lab/.style={inner sep=1pt, align=center, font=\footnotesize}, sm/.style={inner sep=1pt, align=center, font=\scriptsize},
  ttl/.style={anchor=west, inner sep=0pt, font=\small}]
% written by gen/mk_data.py: schematic clouds in whitened coordinates (unit = one ID standard deviation), Delta = 1.768
\def\iddots{1.141/-0.915, 0.793/-0.396, -0.460/0.787, 0.258/-0.211, -0.391/-0.483, 1.610/-0.506, 0.348/-0.784, 1.777/0.323, -2.663/0.210, -1.614/1.601, 0.070/-0.577, -0.432/-0.972, 0.960/-1.231, -0.113/-1.618, -0.096/-0.026, -0.013/-0.091, 1.323/0.224, 0.288/0.800, 1.522/-0.368, -0.714/0.926, -1.364/0.934, 1.369/0.839, 1.282/-2.040, -0.463/-0.773, 0.823/-0.255, 0.763/-0.393, 1.487/-1.352, 1.386/0.398, -0.270/0.410, 0.430/0.463, 0.877/-0.427, -0.372/-0.239, 1.049/-0.658, 0.223/-1.839, -0.944/-0.816, 0.348/0.524, -1.523/0.989, 2.220/-0.857, -1.523/1.908, -0.124/-0.016, -0.895/0.780, -1.154/-0.740, -0.883/-0.517}
\def\ordots{-1.181/-0.324, -2.756/0.911, 0.444/0.636, -3.820/-0.668, -1.408/0.078, -1.856/1.755, -0.855/0.710, -3.603/0.487, -2.943/1.209, -1.415/-1.033, -1.930/0.523, -0.329/-0.267, -2.925/-0.010, -2.026/0.355, -2.910/-0.731, -2.447/-0.342, -1.286/1.351, -3.754/-0.026, -1.556/0.315, -3.228/-0.173, -2.023/0.231, -2.999/-0.910, 0.621/0.860, -1.495/0.729, -2.527/0.524, -0.797/0.393, -0.870/0.235, -1.771/0.334, -2.710/-1.838, -3.078/-1.015, -1.412/0.120, -2.912/0.736, -2.381/2.031, -1.829/0.285, -2.386/0.150, -2.385/1.099, -1.830/0.783, -2.075/-0.558, -1.726/1.890, -2.952/-0.393, -1.908/-0.029, -2.152/-1.155, 0.660/-0.254, -2.828/0.498}
\def\mDelta{1.768}

\def\plane{% frame, clouds and the two means (whitened units)
  \draw[black!35, line width=0.4pt] (-4.4,-2.35) rectangle (3.0,2.35);
  \begin{scope}
    \clip (-4.4,-2.35) rectangle (3.0,2.35);
    \foreach \a/\b in \ordots \fill[black!40] (\a,\b) circle (1.15pt);
    \foreach \a/\b in \iddots \fill[navyc] (\a,\b) circle (1.15pt);
  \end{scope}
  \fill[black!60] (-\mDelta,0) circle (2.1pt);
  \fill[navyc] (0,0) circle (2.1pt);
  \node[lab, anchor=north, fill=white, inner sep=0.5pt] at (-\mDelta,-0.15) {$\vec{\mu}_{o}$};
  \node[lab, anchor=north, text=navyc, fill=white, inner sep=0.5pt] at (0.05,-0.15) {$\vec{\mu}$};
  \draw[navyc, line width=0.9pt] (0.5,0.8) circle (3.0pt);
  \fill[redc] (-3.3,-0.8) circle (2.8pt);}
% ================================================================ (a)
\node[ttl] at (0,2.45) {\textbf{(a)}\ $d^{2}$：$\vec{\mu}$ からの距離の2乗};
\begin{scope}[shift={(3.7,0)}, x=0.8cm, y=0.8cm]
  \begin{scope}
    \clip (-4.4,-2.35) rectangle (3.0,2.35);
    \foreach \r in {1,2,3,4} \draw[navyc, densely dashed, line width=0.55pt] (0,0) circle (\r);
  \end{scope}
  \plane
  \node[sm, text=navyc, anchor=north, fill=white, inner sep=1pt, align=center] at (-0.7,-2.42) {$d^{2}$ が等しい線：$\vec{\mu}$ を中心とする円\\どの向きへ外れても $d^{2}$ は大きくなる};
\end{scope}
% ================================================================ (b)
\node[ttl] at (8.3,2.45) {\textbf{(b)}\ $\ell$：$\vec{\mu}-\vec{\mu}_{o}$ の向きの位置(数1つ)};
\begin{scope}[shift={(11.9,0)}, x=0.8cm, y=0.8cm]
  \begin{scope}
    \clip (-4.4,-2.35) rectangle (3.0,2.35);
    \foreach \x in {-4,-3,...,2} \draw[black!55, densely dashed, line width=0.5pt] (\x,-2.35) -- (\x,2.35);
  \end{scope}
  \plane
  \draw[->, line width=1.0pt] (-\mDelta+0.12,0) -- (-0.12,0);
  \node[sm, anchor=south, fill=white, inner sep=0.5pt] at (-0.88,0.14) {$\vec{\mu}-\vec{\mu}_{o}$};
  \node[sm, anchor=south west, fill=white, inner sep=1pt, align=left] at (-4.35,1.45) {$\ell$ が等しい線：\\$\vec{\mu}-\vec{\mu}_{o}$ に垂直な直線};
  % the number line of l (l = Delta * horizontal coordinate)
  \def\ny{-4.75}
  \draw[->, line width=0.7pt] (-4.4,\ny) -- (3.15,\ny);
  \node[lab, anchor=west] at (3.2,\ny) {$\ell$};
  \draw[black!60, line width=0.5pt, domain=-4.4:3.0, samples=60] plot (\x,{\ny+1.1*exp(-0.5*((\x+\mDelta)/0.93)^2)});
  \draw[navyc, line width=0.8pt, domain=-3.6:3.0, samples=60] plot (\x,{\ny+1.1*exp(-0.5*(\x)^2)});
  \draw[line width=0.6pt] (0,\ny-0.1) -- (0,\ny+0.1);
  \draw[line width=0.6pt] (-\mDelta,\ny-0.1) -- (-\mDelta,\ny+0.1);
  \node[sm, anchor=north] at (0,\ny-0.12) {$0$};
  \node[sm, anchor=north] at (-\mDelta,\ny-0.12) {$-\Delta^{2}$};
  \node[sm, text=navyc, anchor=south west, align=left] at (1.3,\ny+0.55) {分布内の画像\\平均 $0$,標準偏差 ${\approx}\,\Delta$};
  \node[sm, text=black!65, anchor=south east, align=right] at (-3.55,\ny+0.3) {向きを変えた画像\\平均 $-\Delta^{2}$};
  % the two test images dropped onto the line
  \draw[navyc, densely dotted, line width=0.6pt] (0.5,0.68) -- (0.5,\ny);
  \draw[redc, densely dotted, line width=0.6pt] (-3.3,-0.92) -- (-3.3,\ny);
  \fill[navyc] (0.5,\ny) circle (2.2pt);
  \fill[redc] (-3.3,\ny) circle (2.2pt);
\end{scope}
\node[lab, anchor=north west, align=left] at (0.3,-3.0) {(b)の $\ell$ は矢印(ベクトル)ではなく,\\数直線上の1点(スカラー)である.\\画像1枚に $\ell$ が1つ決まる\\(右下の数直線. 横の位置の $\Delta$ 倍).};
\end{tikzpicture}
